\emph{Setting and the imported input.} Work in $\Cl(1,3)$ with $\gamma_{0}^{2} = 1$, $\gamma_{i}^{2} = -1$, $I_{4} = \gamma_{0}\gamma_{1}\gamma_{2}\gamma_{3}$ and $Q(\xi) = \xi_{0}^{2} - \xi_{1}^{2} - \xi_{2}^{2} - \xi_{3}^{2}$. The one input from outside the store is the massless two-point function $D_{+}(\xi) = \int d\mu(k)\, e^{-I_{4}k\cdot\xi}$ over the forward null cone with the invariant measure (numbers, the field theory's phases among them, read in $\mathrm{span}\{1, I_{4}\}$); it is the boundary value $-1/(4\pi^{2}(Q(\xi) - I_{4}0\,\xi_{0}))$, the infinitesimal shift written in scalar form in $\mathrm{span}\{1, I_{4}\}$, so on the open spacelike set $\{Q < 0\}$ it is the real function $-1/(4\pi^{2}Q(\xi))$, even under $\xi \to -\xi$ (standard; Fact \claimref{synthesis.imported-free-field-spin-statistics-and-guido-longo-theorem}(i); e.g. the free-field chapters of Streater and Wightman or Birrell and Davies). Only $f = 1/Q$ and its evenness enter the argument and the checks; the constant $-1/(4\pi^{2})$ and the powers of $I_{4}$ are overall nonzero factors and are kept symbolic. The method is Pauli's free-field argument (Pauli 1940, standard; Fact \claimref{synthesis.imported-free-field-spin-statistics-and-guido-longo-theorem}(ii)); the new content is that the series' own polarisation sums and storeys reproduce it.

\emph{The bracket of a free field.} Let a free field have particle part $\int d\mu\, u(k)a(k)e^{-I_{4}k\cdot x}$ and antiparticle part $\int d\mu\, v(k)b^{*}(k)e^{I_{4}k\cdot x}$, with $[a(k), a^{*}(k')]_{\eta} = [b(k), b^{*}(k')]_{\eta} = \delta$ (the bracket $[x, y]_{\eta} = xy + \eta yx$, $\eta = +1$ for the CAR, $-1$ for the CCR), and suppose the antiparticle sum $\sum vv^{*}$ is $c$ times the particle sum $P(k) = \sum uu^{*}$, $c \neq 0$ real, with $c > 0$ when the Fock space has a positive norm, as for the stored spinor field (Proposition \claimref{field-theory.second-quantisation}(2): the anticommutator is $\sum_{j} 2\alpha_{j}\cos(k_{j}\cdot\xi)\Lambda(k_{j})\chi$, the same $\Lambda(k)$ on both exponentials) and the Hermitian photon (Proposition \claimref{field-theory.fock-space-maxwell-modes}(4)); for a scalar $P = 1$ and positivity of $\langle bb^{*}\rangle$ forces $c > 0$. On an indefinite (Krein) space $\langle bb^{*}\rangle$ may be negative, and then $c < 0$ is allowed. Then the bracket of the field with its adjoint is the c-number
\[
\int d\mu\,P(k)e^{-I_{4}k\cdot\xi} + \eta c\int d\mu\,P(k)e^{I_{4}k\cdot\xi} = P(I_{4}\partial)D_{+}(\xi) + \eta c\,P(-I_{4}\partial)\bigl[D_{+}(-\,\cdot\,)\bigr](\xi),
\]
$\xi = x - y$, where $P$ is any polynomial agreeing with the sum on the cone (two such differ by $Q(k)R(k)$, and $R(I_{4}\partial)\Box D_{+} = 0$ off the cone since $\Box(1/Q) = 0$ there, checked). For $P$ homogeneous of degree $d$, $P(-I_{4}\partial) = (-1)^{d}P(I_{4}\partial)$, and on the spacelike set $D_{+}(-\xi) = D_{+}(\xi)$, so there the bracket is $-\tfrac{I_{4}^{d}}{4\pi^{2}}(1 + \eta c(-1)^{d})P(\partial)f$. The script evaluates the two terms separately, the second from the derivatives of $1/Q$ at $-\xi$ with the chain-rule signs, so the parity is computed, not assumed.

\emph{(1).} Scalar: $P = 1$, $d = 0$. Spinor: by Proposition \claimref{field-theory.second-quantisation}(1) the one-particle projector is $\Lambda(k) = k\gamma_{0}/(2k_{0})$, the factor $1/(2k_{0})$ belonging to the invariant measure, so $P(k) = k\gamma_{0}$, linear: $P(-k) = -P(k)$. The script checks $\Lambda^{2} = \Lambda$, $k\Lambda = 0$, $P(-k) = -P(k)$ and $P(2k) = 2P(k)$ at five rational null $k$. Photon: by Proposition \claimref{field-theory.fock-space-maxwell-modes}(2) the modes at null $k$ are the bivectors $f$ with $kf = 0$; for $\epsilon$ orthogonal to $k$, $k(k \wedge \epsilon) = Q(k)\epsilon - (k\cdot\epsilon)k = 0$, so the modes are $k \wedge \epsilon$ with $\epsilon$ transverse (the script checks two independent transverse $\epsilon$ at each $k$). The polarisation sum on test bivectors is $\sum_{ij}(\delta_{ij} - k_{i}k_{j}/|\mathbf{k}|^{2})\langle A(k \wedge \gamma_{i})\rangle_{0}\langle B(k \wedge \gamma_{j})\rangle_{0}$; the terms of $-g^{\mu\nu}$ beyond the transverse projector are proportional to $k^{\mu}$ or $k^{\nu}$ and vanish under $k \wedge$, so on the cone it equals the quadratic polynomial of (1). The script verifies the equality for all $36$ pairs of basis bivectors at five null $k$, and $P(-k) = P(k)$, $P(2k) = 4P(k)$. So $d = 0, 1, 2 = 2s$ for the three fields.

\emph{(2).} $\partial_{\mu}f = -2\xi_{\mu}/Q^{2}$ and $\partial_{\mu}\partial_{\nu}f = -2g_{\mu\nu}/Q^{2} + 8\xi_{\mu}\xi_{\nu}/Q^{3}$, exact at rational points. For the spinor $P(\partial)f = \partial^{\mu}f\,\gamma_{\mu}\gamma_{0} = -2\xi\gamma_{0}/Q^{2} \neq 0$. For the photon the script computes the coefficients of the quadratic $P(k)[A, B]$ by polarisation and evaluates $P(\partial)f$ for $A = B = \gamma_{0}\gamma_{1}$, nonzero; for the scalar $f \neq 0$. At four rational spacelike points and for $\eta = \pm 1$, $c \in \{1, 2, \tfrac12\}$, the computed bracket vanishes exactly when $c = 1$ and $\eta = -(-1)^{d}$, and with $c = -1$ added to the sample exactly when $\eta = (-1)^{d}$. Since the bracket is $-\tfrac{I_{4}^{d}}{4\pi^{2}}(1 + \eta c(-1)^{d})P(\partial)f$ on the whole spacelike set and $P(\partial)f$ is a nonzero rational function there, the field is local (its bracket vanishes at every spacelike separation, the locality of Proposition \claimref{field-theory.second-quantisation}(2)) if and only if $1 + \eta c(-1)^{d} = 0$, that is, for $c > 0$, $c = 1$ and $\eta = -(-1)^{d}$, and for $c = -1$, $\eta = (-1)^{d}$. The second case needs a negative-norm antiparticle term: a local anticommuting scalar ($d = 0$, $\eta = +1$, $c = -1$) is the Faddeev-Popov ghost of covariant gauge fixing, which lives on a Krein space (\claimref{synthesis.self-adjoint-brst-charge-needs-a-krein-form-and-the-tower-supplies-the-ghosts}); so every statement below about which storey carries a local field holds in a positive-norm space. With $\eta = +1$ and $d = 1$ this is the stored anticommutator $\chi(\nabla D)\gamma_{0}$, a derivative of the Pauli-Jordan function.

\emph{(3).} A tower plane is realised in $\Cl(1,1)$ by $v = \tfrac12(e + f)$, $w = \tfrac12(e - f)$, $e^{2} = 1$, $f^{2} = -1$: $Q(v) = Q(w) = 0$ and $vw + wv = 1$ (the CAR of Fact \claimref{field-theory.imported-tower-limit}), while $vw - wv \neq 1$. The palatial pair on polynomials (Fact \claimref{field-theory.imported-statistics-no-go-palatial-algebra}), $a = d/dz$, $a^{*} = z$, has $[a, a^{*}]z^{n} = z^{n}$ and $\{a, a^{*}\}z^{n} = (2n + 1)z^{n}$, checked for $n \le 5$. So tower modes give $\eta = +1$, palatial modes $\eta = -1$, and by (2), with positive norm, a tower lift is local only for $d$ odd (the spinor) and a palatial lift only for $d$ even (the scalar and the photon). This is the twistor attachment's sentence (Theorem \claimref{field-theory.twistor-attachment}) that a CAR lift is local exactly when $2s$ is odd, now for these attached free fields from locality with positive norm and the sums, the Jost-point theorem being traded for those premises (it stays the input for interacting fields). By Proposition \claimref{field-theory.no-canonical-commutation-relations-tower} no floor and no bounded element of the CAR limit carries the CCR, so a local elementary scalar of positive norm, in particular a quantised rung field read as elementary, is a palatial field; the tower's rung vectors are odd CAR elements and as the modes of a local scalar field would violate (2). Bosons made of tower fermions are not excluded by this (they are even composites, with the CCR reached only in the limit of Proposition \claimref{field-theory.how-close-tower-comes-ccr}).

\emph{Scope.} The result covers the free massless scalar, spinor and Maxwell fields with the stated structure of particle and antiparticle sums; it assumes locality, positive energy and a positive norm, and does not cover interacting fields or the general Wightman setting. For a massive field the same argument needs only that the massive two-point function is real, even and nonzero at spacelike separation (standard; its Bessel-function form is cited from memory and not used). The sign $(-1)^{2s}$ is also the factor by which the boost rotor continued over one modular period ($2\pi$ in imaginary rapidity, Theorem \claimref{synthesis.boost-evolution-on-generators-is-one-unitary-group}(iv)) acts on values; that coincidence is not used here, and a derivation from modular covariance rests on a different premise (\claimref{synthesis.modular-period-on-values-is-the-dual-plane-turn}). Transport-convenient: the sums and storeys are the series' objects; the locality argument itself is standard.
